%%FRAG s:sun
%%BEFORE: \subsubsection*{Khai báo mức module (dòng 22--30)}
\subsubsection{Cơ sở toán: Mặt trời, ngày--đêm và bản đồ bóng}\label{s:sun}
Phần này gồm quỹ đạo mặt trời (\code{Lighting.sun_direction}), máy trạng thái \code{DayCycle} và thiết lập bản đồ bóng (\code{Lighting.shadow_setup}). Phần dùng bản đồ bóng trong shader (\code{shadowFactor}, \code{shadow.vert}) xem mục \ref{s:phong} (\code{splat_prep.vert}) và \code{shadow.vert}. Mô hình Phong và pháp tuyến gaussian cũng nằm ở mục \ref{s:phong}.

\paragraph{Mặt trời (\code{Lighting.sun_direction}).}
\begin{align*}
\mathrm{angle}&=\tfrac{\mathrm{hour}-6}{12}\pi\quad(\text{giờ }6\to0,\;12\to\tfrac\pi2,\;18\to\pi)\\
\mathrm{elevation}&=\sin(\mathrm{angle})\cdot65^\circ\quad(\text{âm ban đêm})\\
\mathrm{azimuth}&=\mathrm{azimuth}_0+(\mathrm{angle}-\tfrac\pi2)\\
\mathrm{dir}&=(\cos el\cos az,\;\sin el,\;\cos el\sin az)\quad(\text{cầu}\to\text{Descartes, đơn vị})
\end{align*}
Chu kỳ: \code{hour += dt*24/day_length (mod 24)}, \code{day_length = 36 s}, tức 1.5 s/giờ.
\begin{itemize}
\item $\mathrm{daylight}=\mathrm{smoothstep}(-0.08,0.3,\mathrm{dir}_y)$.
\item Màu mặt trời: $\mathrm{mix}((1,.5,.25),(1,.96,.88),\mathrm{smoothstep}(0,.35,s))\cdot I\cdot\mathrm{smoothstep}(-.02,.2,s)$ (cam lúc thấp, trắng lúc cao).
\item Ambient: $\mathrm{base}\cdot\mathrm{mix}((.12,.16,.32),1,\mathrm{daylight})$ (xanh tối ban đêm).
\item Bầu trời: $\mathrm{mix}(\text{đêm},\text{ngày},\mathrm{daylight})$ rồi trộn thêm cam với trọng số $0.55\exp(-(s/0.18)^2)$, một bump Gauss quanh đường chân trời (hoàng hôn).
\item $\mathrm{mix}(a,b,t)=a+(b-a)t$.
\end{itemize}

\paragraph{Máy trạng thái ngày--đêm (\code{DayCycle}).}
Mốc bắt đầu: DAWN 5h, DAY 7h, DUSK 17h, NIGHT 19h. \code{phase_for(h)}: $h\in[\mathrm{start},\mathrm{end})$, hoặc nếu $\mathrm{start}>\mathrm{end}$ (qua nửa đêm) thì $h\ge\mathrm{start}\lor h<\mathrm{end}$. Sự kiện \code{on_enter} bật/tắt đèn \code{auto_night}.

\paragraph{Bản đồ bóng (\code{Lighting.shadow_setup}, \code{shadow.vert}, \code{shadowFactor}).}
\begin{itemize}
\item $r=1.1\cdot\mathrm{radius}$ cảnh; mắt đèn $\mathrm{eye}=\mathrm{center}+d\cdot2r$; $\mathrm{up}=(0,1,0)$ (hoặc $(0,0,1)$ nếu $|d_y|\ge0.99$ để tránh $f\parallel\mathrm{up}$).
\item $VP=\mathrm{Ortho}(r,0.01r,4r)\cdot\mathrm{LookAt}(\mathrm{eye},\mathrm{center})$; độ phân giải $2048^2$.
\item $\mathrm{texel}=2r/2048$; $\mathrm{bias}=3\,\mathrm{texel}/(3.99r)$ (độ lệch độ sâu quy về NDC $[0,1]$; miền $z$ dài $3.99r$).
\item Điểm tra: $lp=VP\cdot(p+n\cdot2\,\mathrm{texel})$ (đẩy dọc pháp tuyến chống tự bóng); $s=lp_{xyz}/lp_w\cdot0.5+0.5$ (NDC sang $[0,1]$ để lấy mẫu texture); nằm ngoài $[0,1]$ thì coi là sáng.
\item \textbf{PCF $3\times3$:} $S=\frac19\sum_{i,j=-1}^{1}\big[s_z-\mathrm{bias}\le D(s_{xy}+(i,j)/2048)\big]$.
\item Chỉ gaussian có $\alpha\ge0.3$ và chưa xóa mới được vẽ vào depth map (dạng điểm).
\end{itemize}
%%END

%%FRAG s:phong
%%BEFORE: \subsubsection{Khai báo đầu vào, uniform và biến ra}
\subsubsection{Cơ sở toán: Ánh sáng Phong và bóng đổ}\label{s:phong}
Các hàm \code{shade} và \code{shadowFactor} trong \code{splat_prep.vert} tính ánh sáng cho từng gaussian. Quỹ đạo mặt trời, máy trạng thái ngày--đêm và thiết lập bản đồ bóng (phía CPU, \code{core/lighting.py}) nằm ở mục \ref{s:sun}.

\paragraph{Pháp tuyến gaussian.}
$n$ là trục ngắn nhất của ellipsoid: $\mathrm{axis}=\arg\min|s|$, $n=\mathrm{normalize}(\mathrm{linear}\cdot R[\mathrm{axis}])$ (cột $R$ tương ứng, đưa sang thế giới), lật nếu $n\cdot\mathrm{toCam}<0$. Với gaussian tròn (không có mặt phẳng): $\mathrm{roundness}=\mathrm{smoothstep}(0.35,0.9,s_{\min}/s_{\max})$, $n=\mathrm{normalize}(\mathrm{mix}(n,(0,1,0),\mathrm{roundness}))$. Hàm $\mathrm{smoothstep}(a,b,x)$: $t=\mathrm{clip}((x-a)/(b-a),0,1)$, kết quả $t^2(3-2t)$ (đa thức Hermite bậc 3).

\paragraph{Mô hình Phong.}
\begin{align*}
w(n,l)&=\mathrm{clip}\!\left(\frac{n\cdot l+0.35}{1.35},0,1\right)\quad\text{(wrap diffuse: $n\cdot l=-0.35\mapsto0$, $1\mapsto1$)}\\
I_a&=c\,L_a\\
I_d&=c\sum_j L_j\,w(n,l_j)\,S_j\\
I_s&=k_s\sum_j L_j\max(r_j\cdot v,0)^m S_j,\quad r=\mathrm{reflect}(-l,n)=2(n\cdot l)n-l,\ \text{chỉ khi }n\cdot l>0\\
c_{\mathrm{out}}&=I_a+I_d+I_s\quad\text{(cuối cùng clamp 0..1)}
\end{align*}
$v=\mathrm{normalize}(\mathrm{camPos}-p)$; $m=$ \code{shininess} (mặc định 24); $k_s=$ \code{specular} (0.25).

\paragraph{Đèn điểm: suy giảm.}
$\mathrm{att}=\mathrm{clamp}(1-d^2/\mathrm{range}^2,0,1)^2$ (bình phương nên mượt về 0 tại $d=\mathrm{range}$). Tối đa 8 đèn, không đổ bóng. Màu $=$ màu $\times$ cường độ.
%%END

%%FRAG s:texture
%%BEFORE: \subsubsection*{Khai báo mức module (dòng 14--49)}
\subsubsection{Cơ sở toán: Texture}\label{s:texture}
Các công thức dưới đây giải thích cách dựng UV từ skin, hộp nhân vật, shader avatar (\code{avatar.vert}, \code{avatar.frag}) và bộ giải mã PNG \code{load_png} trong file này.
\begin{itemize}
\item \textbf{Skin $64\times64$} kiểu chữ thập Minecraft, \code{avatar_renderer._uv_region(face, ox, oy, w, h, d)} ($w,h,d$ là rộng, cao, sâu của hộp). Mặt trước $(ox+d,\,oy+d)$; sau $(ox+2d+w,\,oy+d)$; phải $(ox,\,oy+d,\,d\times h)$; trái $(ox+d+w,\,oy+d)$; trên $(ox+d,\,oy)$; dưới $(ox+d+w,\,oy)$. $\mathrm{UV}=\mathrm{pixel}/64$. Bốn đỉnh mỗi mặt: $(r_x,r_y)$, $(r_x+r_w,r_y)$, $(r_x+r_w,r_y+r_h)$, $(r_x,r_y+r_h)$.
\item Hộp: nửa kích thước $h_x=w/2+\mathrm{inflate}$, \ldots (lớp áo ``inflate'' phồng ra để tránh z-fighting). Đỉnh $=$ tâm $+(\pm1)\cdot$nửa kích thước; pháp tuyến là 6 hướng trục.
\item \textbf{Avatar shader:} $\mathrm{shade}=\mathrm{bright}\cdot(0.62+0.38\max(n\cdot l,0))$ (ambient 0.62 + Lambert 0.38); $n_{\mathrm{world}}=\mathrm{normalize}(\mathrm{mat3}(M)\,n)$; bỏ texel $\alpha<0.5$ (lớp áo trong suốt).
\item \textbf{Giải mã PNG tự viết} (\code{load_png}): zlib + bộ lọc dòng, với $a$ là trái, $b$ là trên, $c$ là trên-trái và mọi phép cộng đều mod 256: 1 Sub: $x+a$; 2 Up: $x+b$; 3 Average: $x+\lfloor(a+b)/2\rfloor$; 4 \textbf{Paeth}: $p=a+b-c$, chọn trong $\{a,b,c\}$ gần $p$ nhất ($p_a=|p-a|$, $p_b$, $p_c$; ưu tiên $a\le b\le c$).
\item Texel buffer: dữ liệu gaussian đọc bằng \code{texelFetch} (số nguyên, không nội suy).
\end{itemize}
%%END
